% Compile with: lualatex cv.tex
\documentclass[11pt]{article}

\usepackage[top=0.85in, bottom=0.85in, left=0.9in, right=0.9in]{geometry}
\usepackage{fontspec}
\usepackage{xcolor}
\usepackage{amssymb}
\usepackage{microtype}
\usepackage{fancyhdr}
\usepackage{hyperref}

% ---- Palette (matches niccarp11.github.io light theme) --------------------
\definecolor{accent}{HTML}{B8622C}
\definecolor{accentstrong}{HTML}{94491D}
\definecolor{ink}{HTML}{1C1A17}
\definecolor{muted}{HTML}{5B564E}
\definecolor{faint}{HTML}{8A8478}
\definecolor{rulecol}{HTML}{D8D2C6}

% ---- Fonts (shipped with the site) ----------------------------------------
\setmainfont{spectral}[
  Path = ../fonts/spectral/,
  Extension = .ttf,
  UprightFont = *-regular,
  BoldFont = *-bold,
  ItalicFont = *-italic,
]
\newfontfamily\mono{iosevka-term}[
  Path = ../fonts/iosevka/,
  Extension = .ttf,
  UprightFont = *-regular,
  BoldFont = *-bold,
  ItalicFont = *-italic,
]
\newfontfamily\tracked{iosevka-term}[
  Path = ../fonts/iosevka/,
  Extension = .ttf,
  UprightFont = *-regular,
  BoldFont = *-bold,
  LetterSpace = 6.0,
]

\color{ink}
\linespread{1.04}

\hypersetup{
  colorlinks=true,
  urlcolor=accentstrong,
  linkcolor=accentstrong,
  pdftitle={Nicolò Carpentieri --- Curriculum Vitae},
  pdfauthor={Nicolò Carpentieri},
}

% ---- Section headings: numbered, mono, accented, ruled --------------------
\newcounter{cvsec}
\newcommand{\snum}{\ifnum\value{cvsec}<10 0\fi\arabic{cvsec}}
\newcommand{\cvsection}[1]{%
  \vspace{11pt}\refstepcounter{cvsec}%
  {\tracked\footnotesize\bfseries\textcolor{accent}{\snum}\hspace{9pt}%
   \textcolor{ink}{\MakeUppercase{#1}}}\par%
  \vspace{3.5pt}{\color{rulecol}\hrule height 0.6pt}\vspace{7pt}%
}

% ---- Entry: bold title, right-aligned mono date ---------------------------
\newcommand{\entry}[3]{%
  \textbf{#1}\hfill{\mono\footnotesize\textcolor{muted}{#2}}\\
  #3%
}

% ---- Reusable bits --------------------------------------------------------
\newcommand{\me}{\textbf{Nicolò Carpentieri}}
\newcommand{\venue}[1]{{\mono\footnotesize\textcolor{accentstrong}{#1}}}
\newcommand{\authors}[1]{{\small\textcolor{muted}{#1}}}

% ---- Page furniture -------------------------------------------------------
\pagestyle{fancy}
\fancyhf{}
\fancyhead[L]{\mono\scriptsize\textcolor{faint}{Nicolò Carpentieri}}
\fancyhead[R]{\mono\scriptsize\textcolor{faint}{\thepage}}
\renewcommand{\headrulewidth}{0pt}
\fancypagestyle{firstpage}{\fancyhf{}\renewcommand{\headrulewidth}{0pt}}

\setlength{\parindent}{0pt}

% Tighter, evenly-spaced publication list without enumitem
\newenvironment{cvlist}{%
  \begin{list}{\mono\footnotesize\textcolor{accent}{[\arabic{enumi}]}}{%
    \usecounter{enumi}%
    \setlength{\leftmargin}{2.1em}\setlength{\labelwidth}{1.6em}%
    \setlength{\labelsep}{0.5em}%
    \setlength{\itemsep}{6pt}\setlength{\parsep}{0pt}\setlength{\topsep}{3pt}%
  }%
}{\end{list}}

\begin{document}
\thispagestyle{firstpage}

% ---- Masthead -------------------------------------------------------------
{\mono\fontsize{23}{27}\selectfont\bfseries N\textcolor{accent}{i}colò Carpentieri}\\[3pt]
{\tracked\footnotesize\textcolor{faint}{\MakeUppercase{Systems \& Computer Architecture}}}

\vspace{7pt}
\noindent
\begin{minipage}[t]{0.60\textwidth}
PhD Candidate, Systems Research Group\\
Technical University of Munich
\end{minipage}%
\hfill
\begin{minipage}[t]{0.38\textwidth}
\raggedleft
{\mono\footnotesize
\href{https://niccarp11.github.io}{niccarp11.github.io}\\
\href{mailto:nic.carp11@gmail.com}{nic.carp11@gmail.com}}
\end{minipage}

\vspace{8pt}
{\color{rulecol}\hrule height 0.9pt}

\cvsection{Research Interests}
My interests lie in \textbf{computer architecture and systems}, particularly \textbf{memory disaggregation} at rack scale and the scale-up interconnects, like CXL and related fabrics, that make it possible. Within that space, I focus on \textbf{cache coherence protocols} for heterogeneous accelerators, including GPUs, and on automating the design and verification of the controllers that implement them.

\cvsection{Education}
\entry{Ph.D. in Computer Science}{Oct 2024 -- present}{
Technical University of Munich\\
\textit{Advisor}: Prof. Dr. Pramod Bhatotia
}

\vspace{7pt}
\entry{M.Sc. in Computer Engineering}{2021 -- 2024}{
Politecnico di Torino\\
\textit{Thesis}: Designing and Evaluating Mapping of CNN Layers on an Edge-CGRA, in collaboration with the Embedded Systems Laboratory (ESL), EPFL
}

\vspace{7pt}
\entry{B.Sc. in Electronic Engineering}{2018 -- 2021}{
Politecnico di Torino
}

\cvsection{Research Experience}
\entry{Visiting Student, Embedded Systems Laboratory (ESL), EPFL}{2023 -- 2024}{
Mapping of convolutional neural network algorithms onto CGRA architectures, in collaboration with DAUIN, Politecnico di Torino.
}

\cvsection{Publications}
\begin{cvlist}
\item Decoupling Protocol, Topology, and Policy for Automated Synthesis of Coherence Controllers\\
\authors{\me, Anatole Lefort, David Schall, Julian Pritzi, Pramod Bhatotia}\\
\venue{IEEE/ACM MICRO 2026}

\item vCXLGen: Automated Synthesis and Verification of CXL Bridges for Heterogeneous Architectures\\
\authors{Anatole Lefort, Julian Pritzi, \me, David Schall, Simon Dittrich, Soham Chakraborty, Nicolai Oswald, Pramod Bhatotia}\\
\venue{ACM ASPLOS 2026}\ \ {\textcolor{accent}{$\bigstar$}}\,\textbf{Best Paper Award}

\item C3: CXL Coherence Controllers for Heterogeneous Architectures\\
\authors{Anatole Lefort, David Schall, \me, Julian Pritzi, Nicolai Oswald, Pramod Bhatotia}\\
\venue{IEEE HPCA 2026}

\item Performance Evaluation of Acceleration of Convolutional Layers on OpenEdgeCGRA\\
\authors{\me, Juan Sapriza, Davide Schiavone, Daniele Jahier Pagliari, David Atienza, Maurizio Martina, Alessio Burrello}\\
\venue{ACM CF 2024}
\end{cvlist}

\cvsection{Talks}
ProtoSLICC: Automated Synthesis of Gem5-based Cache Coherence Controllers\\
\authors{gem5 Workshop @ ISCA 2025, Tokyo, Japan} \hfill {\mono\footnotesize\textcolor{muted}{Jun 2025}}

\cvsection{Academic Service}
\textbf{Artifact Evaluation Committee}\hfill{\mono\footnotesize\textcolor{muted}{2026}}\\
\authors{ACM ASPLOS 2026}

\cvsection{Honors and Awards}
Induction in IEEE-HKN (Mu Nu Chapter), Politecnico di Torino chapter \hfill {\mono\footnotesize\textcolor{muted}{Jun 2022}}

\cvsection{Teaching}
Introduction to Software Engineering (EIST) \textit{--- Examiner} \hfill {\mono\footnotesize\textcolor{muted}{2024--2026}}\\
Systems Programming \textit{--- Course Organizer} \hfill {\mono\footnotesize\textcolor{muted}{2026--2027}}

\cvsection{Supervised Theses}
Disaggregated, Vendor-Generic APUs with CXL\\
\authors{Victor Fritz Trost --- B.Sc. Thesis} \hfill {\mono\footnotesize\textcolor{muted}{SS 2025}}

\cvsection{Technical Skills}
{\mono\small C \textcolor{faint}{\textperiodcentered} C++ \textcolor{faint}{\textperiodcentered} Python \textcolor{faint}{\textperiodcentered} VHDL \textcolor{faint}{\textperiodcentered} Verilog}

\end{document}
